% ============================================================
%  MTH636 – Game Theory | Lecture 22
%  Extensive Form Games with Imperfect Information
% ============================================================
\documentclass[12pt, a4paper]{article}

% ============================================================
%  MTH636 -- Game Theory
%  Notes Preamble — Color-Coded, TikZ-Enhanced
% ============================================================

% --- Core packages ---
\usepackage[T1]{fontenc}
\usepackage[utf8]{inputenc}
\usepackage{lmodern}
\usepackage[margin=2.5cm, top=2.8cm, bottom=2.8cm]{geometry}
\usepackage{microtype}

% --- Math ---
\usepackage{amsmath, amssymb, amsthm, mathtools}
\usepackage{bm}        % bold math
\usepackage{cancel}    % strike-through in math

% --- Color & Boxes ---
\usepackage[dvipsnames, svgnames, x11names, table]{xcolor}
\usepackage[most]{tcolorbox}
\tcbuselibrary{breakable, skins, theorems}

% --- TikZ & Diagrams ---
\usepackage{tikz}
\usepackage{tikz-qtree}
\usetikzlibrary{arrows.meta, positioning, shapes.geometric, fit, backgrounds,
                decorations.pathreplacing, calc, matrix, trees}
\usepackage{pgfplots}
\pgfplotsset{compat=1.18}

% --- Tables ---
\usepackage{booktabs}
\usepackage{multirow}
\usepackage{array}
\usepackage{tabularx}
\usepackage{makecell}

% --- Typography & Layout ---
\usepackage{enumitem}
\usepackage{parskip}
\usepackage{fancyhdr}
\usepackage{titlesec}
\usepackage{titletoc}
\usepackage{hyperref}
\usepackage{graphicx}
\usepackage{float}
\usepackage{caption}

% --- Color Palette ---
\definecolor{defblue}{HTML}{1A5276}      % Definition header
\definecolor{defbluebg}{HTML}{D6EAF8}    % Definition background
\definecolor{thmgreen}{HTML}{1E8449}     % Theorem header
\definecolor{thmbg}{HTML}{D5F5E3}        % Theorem background
\definecolor{exorange}{HTML}{B7770D}     % Example header
\definecolor{exbg}{HTML}{FEF9E7}         % Example background
\definecolor{warnred}{HTML}{922B21}      % Warning/Key Insight header
\definecolor{warnbg}{HTML}{FADBD8}       % Warning background
\definecolor{noticegray}{HTML}{2C3E50}   % Note header
\definecolor{noticebg}{HTML}{EBF5FB}     % Note background
\definecolor{taketawayyellow}{HTML}{7D6608} % Key Takeaway header
\definecolor{takeawaybg}{HTML}{FDFEFE}   % Takeaway background
\definecolor{accentpurple}{HTML}{6C3483} % Accent / proof
\definecolor{lightgray}{HTML}{F2F3F4}    % shading

% ============================================================
%  TCOLORBOX ENVIRONMENTS
% ============================================================

% -- Definition (Blue) --
\newtcolorbox{defbox}[2][]{
  enhanced, breakable,
  colback=defbluebg, colframe=defblue, coltitle=white,
  fonttitle=\bfseries\sffamily,
  title={Definition: #2},
  attach boxed title to top left={yshift=-2mm, xshift=4mm},
  boxed title style={colback=defblue, rounded corners},
  arc=4pt, boxrule=1.2pt,
  left=6pt, right=6pt, top=8pt, bottom=6pt,
  #1
}

% -- Theorem / Proposition (Green) --
\newtcolorbox{thmbox}[2][]{
  enhanced, breakable,
  colback=thmbg, colframe=thmgreen, coltitle=white,
  fonttitle=\bfseries\sffamily,
  title={Theorem: #2},
  attach boxed title to top left={yshift=-2mm, xshift=4mm},
  boxed title style={colback=thmgreen, rounded corners},
  arc=4pt, boxrule=1.2pt,
  left=6pt, right=6pt, top=8pt, bottom=6pt,
  #1
}

% -- Proposition (Green variant) --
\newtcolorbox{propbox}[2][]{
  enhanced, breakable,
  colback=thmbg, colframe=thmgreen!70!black, coltitle=white,
  fonttitle=\bfseries\sffamily,
  title={Proposition: #2},
  attach boxed title to top left={yshift=-2mm, xshift=4mm},
  boxed title style={colback=thmgreen!70!black, rounded corners},
  arc=4pt, boxrule=1.2pt,
  left=6pt, right=6pt, top=8pt, bottom=6pt,
  #1
}

% -- Example (Orange/Yellow) --
\newtcolorbox{exbox}[2][]{
  enhanced, breakable,
  colback=exbg, colframe=exorange, coltitle=white,
  fonttitle=\bfseries\sffamily,
  title={Example: #2},
  attach boxed title to top left={yshift=-2mm, xshift=4mm},
  boxed title style={colback=exorange, rounded corners},
  arc=4pt, boxrule=1.2pt,
  left=6pt, right=6pt, top=8pt, bottom=6pt,
  #1
}

% -- Key Insight / Warning (Red) --
\newtcolorbox{warnbox}[2][]{
  enhanced, breakable,
  colback=warnbg, colframe=warnred, coltitle=white,
  fonttitle=\bfseries\sffamily,
  title={Key Insight: #2},
  attach boxed title to top left={yshift=-2mm, xshift=4mm},
  boxed title style={colback=warnred, rounded corners},
  arc=4pt, boxrule=1.2pt,
  left=6pt, right=6pt, top=8pt, bottom=6pt,
  #1
}

% -- Note (Gray-Blue) --
\newtcolorbox{notebox}[1][]{
  enhanced, breakable,
  colback=noticebg, colframe=noticegray!60, coltitle=white,
  fonttitle=\bfseries\sffamily,
  title={Note},
  attach boxed title to top left={yshift=-2mm, xshift=4mm},
  boxed title style={colback=noticegray!60, rounded corners},
  arc=4pt, boxrule=0.8pt,
  left=6pt, right=6pt, top=8pt, bottom=6pt,
  #1
}

% -- Key Takeaway --
\newtcolorbox{takebox}[1][]{
  enhanced, breakable,
  colback=takeawaybg, colframe=taketawayyellow, coltitle=white,
  fonttitle=\bfseries\sffamily,
  title={Key Takeaway},
  attach boxed title to top left={yshift=-2mm, xshift=4mm},
  boxed title style={colback=taketawayyellow, rounded corners},
  arc=4pt, boxrule=1.2pt,
  left=6pt, right=6pt, top=8pt, bottom=6pt,
  #1
}

% -- Intuition --
\newtcolorbox{intubox}[1][]{
  enhanced, breakable,
  colback=accentpurple!8, colframe=accentpurple!60, coltitle=white,
  fonttitle=\bfseries\sffamily,
  title={Intuition},
  attach boxed title to top left={yshift=-2mm, xshift=4mm},
  boxed title style={colback=accentpurple!80, rounded corners},
  arc=4pt, boxrule=0.8pt,
  left=6pt, right=6pt, top=8pt, bottom=6pt,
  #1
}

% ============================================================
%  PAYOFF MATRIX HELPER
%  Usage: inside tabular, use \payoff{a}{b} for (a,b)
% ============================================================
\newcommand{\payoff}[2]{$(#1,\;#2)$}
\newcommand{\nash}[2]{\underline{$#1$},\underline{$#2$}}  % underline NE payoffs
\newcommand{\Nset}{\mathbb{N}}
\newcommand{\Rset}{\mathbb{R}}
\newcommand{\Sset}{\mathcal{S}}
\newcommand{\Ui}{u_i}

% ============================================================
%  SECTION FORMATTING
% ============================================================
\titleformat{\section}
  {\large\bfseries\sffamily\color{defblue}}
  {\thesection}{1em}{}[\titlerule]

\titleformat{\subsection}
  {\normalsize\bfseries\sffamily\color{thmgreen!80!black}}
  {\thesubsection}{1em}{}

\titleformat{\subsubsection}
  {\normalsize\itshape\sffamily\color{exorange!80!black}}
  {\thesubsubsection}{1em}{}

% ============================================================
%  HEADER / FOOTER
% ============================================================
\pagestyle{fancy}
\fancyhf{}
\fancyhead[L]{\sffamily\small\color{defblue} MTH636 -- Game Theory}
\fancyhead[R]{\sffamily\small\color{defblue} \leftmark}
\fancyfoot[C]{\sffamily\small\thepage}
\renewcommand{\headrulewidth}{0.4pt}

% ============================================================
%  HYPERREF SETUP
% ============================================================
\hypersetup{
  colorlinks=true,
  linkcolor=defblue,
  urlcolor=thmgreen,
  citecolor=accentpurple,
  bookmarksnumbered=true,
  pdftitle={MTH636 Game Theory Notes},
  pdfauthor={MTH636 Student Notes}
}

% ============================================================
%  MISC HELPERS
% ============================================================
\newcommand{\BR}{\text{BR}}        % Best Response
\newcommand{\NE}{\text{NE}}        % Nash Equilibrium
\newcommand{\SPNE}{\text{SPNE}}    % Subgame Perfect NE
\newcommand{\Dom}{\text{dom}}      % Dominance
\newcommand{\mixs}[1]{\sigma_{#1}} % Mixed strategy
\newcommand{\pures}[1]{s_{#1}}     % Pure strategy


% ============================================================
\begin{document}

% Title
\begin{center}
  {\LARGE\bfseries\sffamily\color{defblue} MTH636 -- Game Theory}\\[6pt]
  {\Large\sffamily Lecture 22: Extensive Form Games with Imperfect Information}\\[4pt]
\end{center}
\vspace{0.5cm}
\hrule height 1.5pt
\vspace{0.3cm}

% ============================================================
\section{Motivation: Imperfect Information in Extensive Form}

So far we have studied extensive form games with \emph{perfect information}, where each player observes the full history of play before making a decision. In many real-world situations, however, a player may not know which node in the game tree they currently occupy. We now introduce the framework of \textbf{extensive form games with imperfect information} (and possible chance moves).

\subsection{Illustrative Examples}

\paragraph{Battle of the Sexes -- two versions.}
Consider the standard Battle of the Sexes scenario with payoffs $(2,1)$ for $(F,F)$, $(0,0)$ for $(F,D)$, $(0,0)$ for $(D,F)$, and $(1,2)$ for $(D,D)$.

\begin{itemize}
  \item \textbf{Left game (Man moves first):} The root is labelled $M$ (Man). Man chooses $F$ or $D$. Woman (W) observes Man's choice and responds. Since Woman can see Man's move, the two Woman nodes are in \emph{separate} information sets -- this is a game of \textbf{perfect information}.
  \item \textbf{Right game (Woman moves first):} The root is labelled $W$ (Woman). Woman chooses $F$ or $D$. Man (M) does \emph{not} observe Woman's choice, so his two decision nodes are connected by a dashed oval, indicating a single information set -- this is a game of \textbf{imperfect information}.
\end{itemize}

In both cases the terminal payoffs are $(2,1)$, $(0,0)$, $(0,0)$, $(1,2)$ at the four leaves.

\paragraph{A richer example with Nature.}
Consider a game tree rooted at $n_0$ (Player I). Player I chooses $A$, $B$, or $C$.
\begin{itemize}
  \item After $A$: reach $n_1$. After $B$: reach $n_2$. Player II has information set $\{n_1, n_2\}$ (dashed oval -- she cannot distinguish $n_1$ from $n_2$). From this information set she plays $\ell$ or $g$. Outcomes: $(4,1)$, $(7,2)$ from $n_1$, and $(6,3)$, $(4,0)$ from $n_2$.
  \item After $C$: reach $n_3$. Nature (O) moves, playing $\tfrac{1}{4}$ toward $n_4$ and $\tfrac{3}{4}$ toward $n_5$. At $n_4$, Player II chooses $L$ or $R$. At $n_5$, Player I chooses $a$ or $b$.
  \item There is also an information set $\{n_1, n_2, n_5\}$ visible in the diagram.
\end{itemize}

% ============================================================
\section{Formal Definition}

\begin{defbox}{Extensive Form Game with Imperfect Information}
A game in extensive form (with imperfect information and chance moves) is a tuple
\[
  \Gamma = \bigl\langle N,\, V,\, E,\, n_0,\;
           (V_i)_{i \in N \cup \{0\}},\;
           \mathcal{H} = (H_i)_{i \in N},\;
           \mathcal{A} = (A(h))_{h \in \bigcup_i H_i},\;
           (p_n)_{n \in V_0},\;
           (u_i)_{i \in N}
           \bigr\rangle
\]
where the components are defined as follows.
\end{defbox}

\subsection{Components of the Tuple}

\begin{itemize}
  \item $N$: the set of \textbf{players}.

  \item $(V, E, n_0)$: a \textbf{directed rooted tree}.

  \item $(V_i)_{i \in N \cup \{0\}}$: a \textbf{partition} of the set of non-leaf nodes among the agents and nature. $V_0$ is the set of \emph{nature's nodes} (chance nodes); $V_i$ is the set of nodes at which player $i$ moves.

  \item $\mathcal{H} = (H_i)_{i \in N}$: $H_i$ is a \textbf{partition} of $V_i$. Each set in the partition is called an \textbf{information set}. Each $h \in H_i$ must satisfy:
    \begin{enumerate}
      \item[(i)] Each node in $h$ has the \textbf{same number of outgoing edges}.
      \item[(ii)] Every path in $(V, E, n_0)$ intersects $h$ \textbf{at most once}.
    \end{enumerate}

  \item $\mathcal{A} = (A(h))_{h \in \bigcup_{i=1}^{n} H_i}$: for each information set $h \in \bigcup_{i} H_i$, $A(h)$ is a \textbf{partition of the edges} coming out from the nodes in $h$, such that for all $a \in A(h)$, $a$ contains exactly one edge outgoing from $n$ for all $n \in h$.

    A set $a \in A(h)$ is called an \textbf{action at} $h$.

  \item $(p_n)_{n \in V_0}$: $p_n$ denotes a \textbf{probability distribution} over the edges outgoing from $n$ (for each chance node $n$).

  \item $(u_i)_{i \in N}$: the \textbf{utility functions} defined at the leaf nodes.
\end{itemize}

% ============================================================
\section{Information Sets: Validity Conditions}

\subsection{Invalid vs.\ Valid Game Trees}

The two conditions on information sets -- (i) equal number of outgoing edges at every node in $h$, and (ii) every root-to-leaf path crosses $h$ at most once -- are essential for the game to be well-defined.

\paragraph{Example of an invalid game.}
Consider a two-player tree: Player I at $n_0$ chooses $L$ or $R$, leading to $n_1$ or $n_2$. From $n_1$, Player II moves and then the game continues to a subtree containing $n_3$ and $n_4$. From $n_2$, the game continues to $n_4$. If $n_3$ and $n_4$ are in the same information set of Player II, then there exists a path from the root that passes through this information set \emph{more than once}, violating condition (ii). This is \textbf{not a valid extensive form game}.

\begin{notebox}
The annotated note in the lecture reads: ``Not a valid extensive form game.''
\end{notebox}

\paragraph{Example of a valid game (not satisfying perfect recall).}
Consider Player I at $n_0$ choosing $L$ or $R$ to reach $n_1$ or $n_2$. Player II has information set $\{n_1, n_2\}$ (she cannot distinguish them). From each she plays $a$, $b$, $d$ (or similar). A second Player II information set groups some descendants together. Every path from the root crosses each information set at most once, so this is a \textbf{valid} extensive form game -- however it \textbf{does not satisfy perfect recall} (noted in red in the lecture).

% ============================================================
\section{Perfect Recall}

\begin{defbox}{Perfect Recall}
An extensive form game satisfies \textbf{perfect recall} for player $i$ if for all information sets $h \in H_i$ and for all $u, y \in h$, an action of player $i$ is on the path from the root to $u$ if and only if it is on the path from the root to $y$.
\end{defbox}

In other words, perfect recall requires that whenever two nodes belong to the same information set of player $i$, the player must have taken exactly the same sequence of actions to reach either node -- she has not ``forgotten'' any earlier choices she made.

\subsection{Example of a Violation of Perfect Recall}

Consider the following game tree:
\begin{itemize}
  \item Player I ($n_0$) chooses $L$ or $R$.
  \item After $L$: Player II ($n_1$) chooses $\ell$ or $r$, reaching $n_3$ or $n_4$.
  \item After $R$: reach $n_2$, then Player II ($n_2$) chooses $\ell$ or $r$, reaching $n_4$ or $n_5$.
  \item Nodes $n_3$, $n_4$, $n_5$ are grouped into a single Player II information set.
  \item Player II chooses $\alpha$ or $\beta$ from this information set.
\end{itemize}

In this game, the path to $n_4$ can arise via two different routes: $n_0 \xrightarrow{L} n_1 \xrightarrow{r} n_4$ (Player II moved at $n_1$) or $n_0 \xrightarrow{R} n_2 \xrightarrow{\ell} n_4$ (Player II moved at $n_2$). In both cases Player II ends up in the information set $\{n_3, n_4, n_5\}$, but whether she earlier chose $r$ (at $n_1$) or $\ell$ (at $n_2$) differs between the two paths. Thus Player II has ``forgotten'' her own earlier move.

\begin{warnbox}{Violation of Perfect Recall}
This game \textbf{violates perfect recall} for Player II (noted in red in the lecture). Perfect recall is violated whenever a player cannot remember their own past actions when entering an information set.
\end{warnbox}

\begin{notebox}
Perfect recall is a standard assumption in most game-theoretic models. When it holds for all players, Kuhn's theorem guarantees that mixed strategies and behaviour strategies are \emph{outcome-equivalent}.
\end{notebox}

\end{document}
